\honest{Categorizing Has Consequences}{Eliezer Yudkowsky}{2008}

Your blood type is real. Karl Landsteiner discovered blood types, which made transfusion safe, and won the 1930 Nobel Prize for it. A mother whose type lacks the ``+'' antigen can react against a second child who has it. ``Thus people learn their blood types before they marry.''

Also, in Japan, type A people are ``earnest and creative'' and type B people ``wild and cheerful,'' and there are blood type horoscopes in the daily paper. This is odd, I say, because blood types were never mysterious, and there has never been a war, or even a political conflict, between them. So ``The stereotypes must have arisen strictly from the mere existence of the labels.'' \nb{As Wikipedia summarizes the history, the stereotypes have authors: Takeji Furukawa's 1927 paper on blood type and temperament, a later study by Furukawa that explained Taiwanese resistance to Japanese rule by the share of type O blood, and the books of Masahiko Nomi in the 1970s. The case shows people believing claims written about a category, not labels acting alone.}

Would the same happen with plants, rocks or office furniture? ``I can't recall reading about such an experiment,'' I say, but I predict that ``mere labeling had power over all things, at least in the human imagination.'' \nb{Such experiments had been done. Carmichael, Hogan and Walter (1932) found that naming a drawing could change how people later reproduced it. Tajfel and Wilkes (1963) found that lines labeled by length were judged farther apart, and also that a classification ``not coherently related'' to length had no such effect.}

I give four ways to see it. A boundary makes the mind hunt for similarities, and ``a weakly negative correlation can be mistaken for a strong positive one with a bit of selective memory.'' A name sets aside a neural subnetwork. Things with one name are compressed into one point. And once you name a category, you can make things up about it. \nb{The first effect has a name in psychology, illusory correlation. The fourth is the one the history of the blood-type theory fits.} It is not just Japan: \textsc{Eat Right 4 Your Type} was a bestseller in the West.

So drawing a boundary is ``not a neutral act.'' Perhaps a purer Bayesian AI could consider a class without being swayed by it; you cannot. That is ``One more reason not to believe you can define a word any way you like.''
